\section{引言}

\subsection{编写目的}
本文档用于定义类 Discord 实时社区交流系统的业务范围、功能需求、接口与数据要求、非功能指标以及验收标准，为课程设计、原型开发、测试设计和后续演示提供统一依据。本文档强调需求层面的完整性与可验证性，不直接展开数据库实现、部署脚本或具体技术选型细节。

\subsection{项目背景}
随着在线学习、兴趣社区协作和远程团队沟通需求持续增长，用户需要一种兼具群组交流、即时消息、语音互动和权限管理能力的实时社区平台。类 Discord 实时社区交流系统面向课程设计场景，聚焦“社区 + 频道 + 实时互动”的核心体验，支持用户围绕不同主题组织文本与语音交流。

\subsection{文档范围}
本软件需求规格说明书覆盖以下内容：
\begin{itemize}
  \item 平台账号体系，包括注册、登录、身份认证、密码找回和个人资料维护。
  \item 社交协作能力，包括好友关系、私聊会话、在线状态与通知提醒。
  \item 社区组织能力，包括社区创建、加入、退出、成员管理、文本频道与语音频道。
  \item 权限控制能力，包括角色定义、频道访问限制和操作授权校验。
  \item 支撑上述业务的接口约束、数据实体定义、性能指标和验收条件。
\end{itemize}

本文档明确不包含以下内容：
\begin{itemize}
  \item 真实多人语音媒体流、屏幕共享和直播能力；v1 仅描述语音频道成员状态同步。
  \item 短视频、机器人平台、第三方开放接口市场等扩展生态能力。
  \item 商业化订阅、广告投放、付费会员和独立运营后台系统。
  \item 具体代码实现、底层数据库建模脚本、自动化部署流程和详细测试用例实现。
\end{itemize}

\subsection{预期读者}
本文档的主要读者包括：
\begin{itemize}
  \item 项目开发人员：依据需求拆分模块、设计接口与实现功能。
  \item 项目测试人员：基于功能需求和验收条件制定测试方案。
  \item 任课教师或评审人员：检查系统目标、需求完整性和课程成果质量。
  \item 项目成员：统一业务术语、边界理解和交付预期。
\end{itemize}

\subsection{系统边界}
本系统围绕实时社区交流场景展开，其边界如下：
\begin{itemize}
  \item 系统内部：用户账号、好友关系、社区空间、频道组织、消息发送、语音状态、角色权限、通知提醒。
  \item 系统外部：邮件或短信验证码服务、文件存储服务、推送通知服务、浏览器与客户端运行环境。
  \item 边界规则：凡涉及社区交流主流程的数据生产、分发、展示与访问控制，均属于系统职责；依赖的基础通信或存储能力由外部服务提供。
\end{itemize}

\subsection{术语与缩略语}
\noindent\begin{tabularx}{\textwidth}{L{3cm}Y}
  \toprule
  术语或缩写 & 说明 \\
  \midrule
  社区（Server） & 用户围绕某一主题或组织建立的交流空间，承载频道、成员、角色和权限规则。 \\
  频道（Channel） & 社区内部的交流单元，分为文本频道和语音频道两类。 \\
  文本频道 & 用于发送文字、图片、文件和提及消息的频道。 \\
  语音频道 & 用于实时语音接入、状态同步和在线互动的频道。 \\
  私聊（Direct Message, DM） & 两个用户之间的一对一会话，不依附于某个社区。 \\
  在线状态（Presence） & 用户当前的可用状态，如在线、离开、忙碌、隐身和离线。 \\
  提及（Mention） & 在消息中使用特定标记引用用户，使其收到重点提醒。 \\
  角色（Role） & 社区内可分配给成员的权限集合，用于控制频道访问和管理操作。 \\
  实时同步 & 客户端与服务端通过长连接或事件机制及时传播状态变化的能力。 \\
  SRS & Software Requirements Specification，软件需求规格说明书。 \\
  \bottomrule
\end{tabularx}

\subsection{参考依据}
\begin{itemize}
  \item 软件工程课程设计任务要求与课程文档规范。
  \item 软件需求规格说明书常见写作规范与高校课程报告格式要求。
  \item 即时通信与在线社区平台的通用业务模式和交互需求。
  \item 面向实时系统的性能、安全、可用性与权限控制基本原则。
\end{itemize}
